\section{Initialization, fitting, and a signed perturbation estimate}
\label{cp:fitting-section}

\begin{lemma}[Initialization and real fitting]
\label{cp:fit}
For this lemma put
\[
 H=\max_{0\le j\le L}\sqrt{Q^{(j)}_{11}},\qquad
 \mathsf H(t)=[h^{(L)}(t,v_1),\ldots,h^{(L)}(t,v_m)].
\]
For the parameter tuple
$\theta=(W^{(1)},\ldots,W^{(L)},w)$ use the norm
\[
\|\theta\|_{\mathrm{par}}^2
 =\|W^{(1)}\|_F^2/n+\sum_{j=2}^L\|W^{(j)}\|_F^2+\|w\|^2/n.
\]
The initialized training Grams satisfy
\begin{equation}
 \left(\frac{h_0^{(j)}(v_a)^\top h_0^{(j)}(v_b)}n\right)_{a,b=1}^m
 \xrightarrow{\mathbb P}Q^{(j)}
 \qquad(1\le j\le L).
 \label{cp:eq-initial-gram-limit}
\end{equation}
With probability tending to one,
\begin{equation}
 \begin{gathered}
 \frac{\|W_0^{(1)}\|_{\mathrm{op}}}{\sqrt n}\le8,\qquad
 \max_{2\le j\le L}\|W_0^{(j)}\|_{\mathrm{op}}\le8,\\
 \sup_{\|v\|=1,\,1\le j\le L}\frac{\|h_0^{(j)}(v)\|}{\sqrt n}
 \le\frac{3H}{2},\qquad
 \frac{\mathsf H_0^\top\mathsf H_0}{mn}\succeq\frac{\gamma}{2m}I_m.
 \end{gathered}
 \label{cp:eq-initialization-event}
\end{equation}
On this event the flow is global, converges in parameters, fits the labels,
and satisfies, for all $t\ge0$,
\begin{gather}
 \begin{gathered}
 \frac{\|W^{(1)}(t)\|_{\mathrm{op}}}{\sqrt n}<9,\qquad
 \max_{2\le j\le L}\|W^{(j)}(t)\|_{\mathrm{op}}<9,\\
 \sup_{\|v\|=1,\,1\le j\le L}\frac{\|h^{(j)}(t,v)\|}{\sqrt n}<2H,\qquad
 \frac{\mathsf H(t)^\top\mathsf H(t)}{mn}\succeq\frac{\gamma}{4m}I_m
 \end{gathered}
 \label{cp:eq-real-tube}\\
 \begin{gathered}
 \frac{\|r(t)\|_2}{\sqrt m}\le Ye^{-\gamma t/(2m)},\qquad
 \frac{\|w(t)\|}{\sqrt n}\le2Y\sqrt{\frac m\gamma},\\
 \int_t^\infty\|\dot\theta(u)\|_{\mathrm{par}}\,du
       \le\frac{2\|r(t)\|_2}{\sqrt\gamma}.
 \end{gathered}
 \label{cp:eq-dense-tail}
\end{gather}
The output converges uniformly on the sphere, with
\begin{equation}
 \sup_{\|v\|=1}|f_n(\infty,v)-f_n(t,v)|
 \le C\frac m\gamma\frac{\|r(t)\|_2}{\sqrt m}.
 \label{cp:eq-dense-output-tail}
\end{equation}
These assertions hold jointly for all rectangular cavities deleting any
fixed bounded number of coordinates per layer, with the original width-$n$
variances, flow factors, and denominators unchanged.
\end{lemma}

\begin{proof}
All auxiliary coefficients in this proof are local. Put
$\lambda=\gamma/m$, $\rho=\|r\|_2/\sqrt m$,
\[
 s=\max\left\{1,\max_j\sup_{|\operatorname{Im}z|\le a/2}
 |\phi_j'(z)|\right\},\qquad b=\max_j|\phi_j(0)|,
\]
and retain $H$ as defined in this lemma's statement. Thus
$|\phi_j(z)|\le b+s|z|$ on the half strip. Define
\begin{align}
 D_j&=s(9s)^{L-j},&
 U_1&=D_1,& U_j&=2HD_j\quad(j\ge2),\nonumber\\
 F_1&=sU_1,&F_j&=s(2HU_j+9F_{j-1}),&F&=F_L.
 \label{cp:eq-fitting-coefficients}
\end{align}
Equal input norms give $Q^{(j)}_{aa}=Q^{(j)}_{11}$, hence
$\lambda\le H^2$. The coefficients satisfy
\begin{equation}
 H\le(2\beta)^L,\qquad
 F=s^2\left[(9s)^{2L-2}+4H^2\sum_{j=0}^{L-2}(9s)^{2j}\right],
 \qquad8H\sqrt F\le\beta^{5L}.
 \label{cp:eq-fitting-powers}
\end{equation}
The first inequality follows from Minkowski's inequality and
$\sqrt{Q^{(j)}_{11}}\le b+s\sqrt{Q^{(j-1)}_{11}}$.
For the last, summing the geometric series gives
$F\le(21/20)H^2s^2(9s)^{2L-2}$, and then
$8H\sqrt F\le H^2(9s)^L\le36^L\beta^{3L}\le\beta^{5L}$.
Consequently the assumed cap
$Y\le\lambda\beta^{-30L}$ implies
$Y\le\lambda/(8H\sqrt F)$, which will close the real-flow estimates.

A maximal $1/4$-separated subset of the unit sphere in $\mathbb R^k$
has at most $9^k$ points: disjoint radius-$1/8$ balls fit inside the
radius-$9/8$ ball.  Approximating each side of a maximizing bilinear
pair gives $\|M\|_{\mathrm{op}}\le2\max|u^\top Mv|$ over these nets.
Each pairing for $W_0^{(\ell)}$, $\ell\ge2$, or
$W_0^{(1)}/\sqrt n$ is $N(0,1/n)$.  The exponential Gaussian tail and
a union bound prove the operator assertions, with failure bounded by
$2(L-1)e^{-(8-2\log9)n}+2e^{-(8-\log9)n+d\log9}$.

On any fixed finite query list, condition on the preceding initialized
layers.  Each next empirical Gram is an average of independent
Gaussian-row feature products, with conditional mean
$\Psi_\ell(C)_{ab}=\mathbb E\phi_\ell(Z_a)\phi_\ell(Z_b)$,
$Z\sim N(0,C)$.  On bounded covariance sets the conditional variances
are bounded, by linear growth and Gaussian fourth moments.  Moreover
$\Psi_\ell$ is continuous on the entire positive semidefinite cone:
write $Z=C^{1/2}G$ with a standard Gaussian vector $G$ of the matching
dimension, use continuity of the positive square root, and
dominate by a fixed Gaussian polynomial.  Conditional Chebyshev and
induction therefore prove convergence of each empirical Gram to
$Q^{(\ell)}$, including when intermediate covariances are singular.
Apply this to the training inputs and a fixed sphere net of mesh
$h\le H/[4(8s)^L]$.  On the operator event the initialized layer-$\ell$
feature map, in neuron RMS, is $(8s)^\ell$-Lipschitz.  Its norm at net
points is eventually at most $\sqrt{5/4}H$, with probability tending to
one; the off-net increment is at most $H/4$.  This proves the sphere
bound and the training Gram assertion in
\eqref{cp:eq-initialization-event}.

Here is also a uniform justification for the cavity assertion, without
a union bound over deletion sets.  On a sphere net of mesh $n^{-2}$,
the preceding RMS bounds and conditional Gaussian tails give
$\max_{i,\ell,v}|h_{0,i}^{(\ell)}(v)|=O(\sqrt{\log(en)})$
with probability tending to one.  The net has only a fixed power of $n$
points; the coordinate Lipschitz bound is at most
$\sqrt n(8s)^\ell$, so the same bound holds off the net.  Zero-extend
each cavity's features into $\mathbb R^n$.  At a deleted layer, at most
a fixed number of original coordinates are removed.  At an undeleted
coordinate the activation difference is at most $s$ times the
preactivation difference.  Thus successive RMS differences are bounded
by $8s$ times the preceding difference plus
$O(\sqrt{\log(en)/n})$.  At fixed depth their supremum, simultaneously
over all such cavities and sphere queries, tends to zero.  Submatrix
operator norms do not increase.  Taking strictly smaller initial
feature and Gram margins before this comparison gives
\eqref{cp:eq-initialization-event} for every cavity.  This also shows
directly that each cavity's training Gram has the same limit $Q^{(L)}$.

Stop any one of these flows at the first failure of
\eqref{cp:eq-real-tube}.  Backward propagation gives
$\|\delta_a^{(\ell)}\|/\sqrt n\le D_\ell\|w\|/\sqrt n$.
In Euclidean mobility coordinates
$(W^{(1)}/\sqrt n,W^{(2)},\ldots,W^{(L)},w/\sqrt n)$,
the vector field is minus the gradient of $\rho^2$.  Consequently
\[
 -\partial_t\rho^2=\|\dot\theta\|_{\mathrm{par}}^2
 \ge\lambda\rho^2.
\]
The last inequality uses only the readout contribution and the stopped
Gram gap.  For any stopped interval $[t,T]$, Cauchy--Schwarz gives
\[
 \left(\int_t^T\|\dot\theta\|_{\mathrm{par}}\right)^2
 \le\left(\int_t^T\frac{\|\dot\theta\|_{\mathrm{par}}^2}{\rho}\right)
      \left(\int_t^T\rho\right)
 \le2\rho(t)\,\frac{2\rho(t)}\lambda.
\]
If a residual vanishes, all velocities vanish, so the estimate extends
without division by zero.  In particular $\|w\|/\sqrt n\le2Y/\sqrt\lambda$.
Sample Cauchy--Schwarz bounds each normalized hidden-block velocity by
$2\rho U_\ell\|w\|/\sqrt n$.  Integration and forward subtraction give
\begin{gather}
 \frac{\|W^{(1)}-W_0^{(1)}\|_F}{\sqrt n}
       \le\frac{8U_1Y^2}{\lambda^{3/2}},\qquad
 \|W^{(\ell)}-W_0^{(\ell)}\|_F
       \le\frac{8U_\ell Y^2}{\lambda^{3/2}},\nonumber\\
 \sup_v\frac{\|h^{(\ell)}-h_0^{(\ell)}\|}{\sqrt n}
       \le\frac{8F_\ell Y^2}{\lambda^{3/2}}.
 \label{cp:eq-dense-displacements}
\end{gather}
Indeed each later feature difference is bounded by $s$ times
$2H\|W^{(\ell)}-W_0^{(\ell)}\|_F$ plus nine times the preceding
feature difference, which is exactly the recursion defining $F_\ell$.
The label restriction implies
$8FY^2/\lambda^{3/2}\le\sqrt\lambda/(8H^2)\le1/(8H)$.
Thus the operators remain below $8+1/8$, the features below
$3H/2+H/8$, and the smallest singular value of
$\mathsf H/\sqrt{mn}$ above
$(1/\sqrt2-1/8)\sqrt\lambda>\sqrt\lambda/2$.
Every stop has a strict margin.  At fixed width the bounded parameters
and locally Lipschitz vector field give global continuation.  The
finite parameter length proves convergence, and residual decay proves
fitting.  The argument uses only the retained dimensions and the
original denominators $n$, and hence also proves cavity fitting.

Finally the same forward recursion gives
\[
 \|\dot w\|/\sqrt n\le4H\rho,\qquad
 \sup_v\|\dot h^{(L)}(v)\|/\sqrt n\le2F\rho\|w\|/\sqrt n.
\]
Therefore
\[
 \sup_{\|v\|=1}|f_n(\infty,v)-f_n(t,v)|
 \le\frac{16}{\lambda}
       \left(H^2+\frac{FY^2}{\lambda}\right)\rho(t).
\]
The coefficients $H,F$ depend only on activations and depth, and the
headline cap with $\lambda\le1$ gives $Y^2/\lambda\le1$.
This proves \eqref{cp:eq-dense-output-tail}, including its uniform endpoint.
\end{proof}

\begin{lemma}[Signed perturbation estimate]
\label{cp:signed}
Let a differentiable prediction map $F$ take Euclidean parameter space
to sample space with inner product
$\langle u,v\rangle_m=m^{-1}u^\top v$.  Write $r=F(\theta)-y$,
$r'=F(\theta')-y$, $\rho=\|r\|_m$, $\rho'=\|r'\|_m$, and
$J=DF(\theta)$, $J'=DF(\theta')$.  Suppose
$\dot\theta=-2J^*r+\mathcal E$ and
$\dot\theta'=-2J'^*r'$. Here $J^*$ and $J'^*$ are adjoints for the
stated parameter and sample inner products. If a constant $K\ge0$
satisfies, for $e=\|\theta-\theta'\|$,
\[
 \|r-r'-J'(\theta-\theta')\|_m\le Ke^2/2,
 \qquad\|(J-J')(\theta-\theta')\|_m\le Ke^2,
\]
then, on every finite interval,
\begin{equation}
 e(t)\le
 \left(e(0)+\int_0^t\|\mathcal E(u)\|\,du\right)
 \exp\left\{K\int_0^t(\rho'+3\rho)\,du\right\}.
 \label{cp:eq-signed-stability}
\end{equation}
\end{lemma}
\begin{proof}
Put $d=\theta-\theta'$ and
$R=r-r'-J'd$.  Exact subtraction gives
\[
 \tfrac12\partial_t e^2
 =-2\|r-r'\|_m^2+2\langle r-r',R\rangle_m
   -2\langle r,(J-J')d\rangle_m+\langle d,\mathcal E\rangle
 \le K(\rho'+3\rho)e^2+e\|\mathcal E\|.
\]
Apply the integrating factor to $\sqrt{e^2+\epsilon^2}$ and let
$\epsilon\downarrow0$.  This proves the assertion also through zeros
of the parameter discrepancy.
\end{proof}
